\section{问题一的模型建立与求解}

\subsection{求解思路}

问题一分三步：先把 22 个指标聚合为记录级与域级质量评分，并在扩展集上对照；再在记录与
指标对两个层面度量质量冲突；最后在单纯形约束下建立配比–损失响应面，按检验集的不同角色
分别验证，并讨论向更大规模外推的可行范围。

\subsection{质量评价模型}

设六个打分体系为 $f=1,\dots,6$（DSIR、RedPajama 规则、FineWeb-Edu、WanJuan、QuRating、
PRRC），体系 $f$ 含指标集合 $J_f$。记录 $i$ 的质量评分取\textbf{体系平衡平均}
\begin{equation}
  Q_i=\frac{1}{6}\sum_{f=1}^{6}\frac{1}{|J_f|}\sum_{j\in J_f}x_{ij},
  \label{eq:q-record}
\end{equation}
即先在体系内平均、再在体系间等权平均。与 22 个指标直接等权相比，式 \eqref{eq:q-record}
使含 11 个指标的规则体系与只含 1 个指标的体系具有相同权重，避免指标数量转化为权重。
域级评分 $Q_d$ 取该域全部记录的平均，总体评分按各域的观测记录数加权。由于秩分在 A1 全体上
校准，A1 总体评分恒为 0.5，这是校准的结果而非独立的质量证据。

域级评分的区间采用域内自助法（1000 次）。A1 中每个域只来自一个源数据分片，因此采用以
文档为单位的条件自助法；区间以已冻结的归一化和评分规则为条件，不包含方法选择带来的不确定性，
也不包含分片内部未能识别的相关性。

\begin{table}[htbp]
  \centering
  \caption{域级质量评分（体系平衡平均）}
  \label{tab:q1-domain}
  \begin{tabular}{lrrrl}
    \toprule
    领域 & 记录数 & $Q_d$ & 95\% 区间 & 扩展集对照 \\
    \midrule
    arXiv & 1,419 & 0.6594 & [0.6566, 0.6621] & A2（17,523 条）：0.6578 \\
    StackExchange & 10,000 & 0.5396 & [0.5379, 0.5414] & \\
    C4 & 10,000 & 0.5165 & [0.5137, 0.5191] & \\
    CommonCrawl & 9,636 & 0.5155 & [0.5133, 0.5177] & \\
    Wikipedia & 9,986 & 0.4863 & [0.4841, 0.4885] & \\
    Book & 171 & 0.4681 & [0.4533, 0.4821] & \\
    GitHub & 10,000 & 0.4206 & [0.4190, 0.4224] & A3（203,751 条）：0.4211 \\
    \bottomrule
  \end{tabular}
\end{table}

表\ref{tab:q1-domain}给出结果。arXiv 最高、GitHub 最低。扩展集 A2（全部为 arXiv）与
A3（全部为 GitHub）与抽样集的差值分别为 $-0.0016$（95\% 区间 $[-0.0043, 0.0012]$）与
$0.0005$（$[-0.0011, 0.0020]$），区间都包含 0，域排序不变。需要说明，扩展集与抽样集
来自同一来源且相互重叠，这一对照检验的是评分对样本量的稳定性，而不是独立来源的重复。

评分对方法选择的敏感性通过 11 种预先声明的变体检验（换用等权或中位数聚合、依次去掉一个
打分体系、改变多维指标的压缩方式与非单调指标的处理）。在全部变体中 arXiv 始终最高、
GitHub 始终最低，中间各域的排序随方法变化。缺失记录按假设 2 取三种极端值时，任一域评分
的变化不超过 0.0002，排序不变。

\subsection{质量冲突的定义、度量与消解}

\textbf{定义。}对记录 $i$，定义其指标分歧幅度
\begin{equation}
  R_i=\max_{j}x_{ij}-\min_{j}x_{ij},
\end{equation}
称 $R_i\ge\tau$ 的记录在阈值 $\tau$ 下存在冲突；在预先固定的阈值网格 $\tau\in\{0.1,\dots,0.9\}$
上报告冲突率曲线 $\Pr(R_i\ge\tau)$，并同时报告记录内指标的平均标准差。在指标对层面，
把 Spearman 秩相关不高于 $-0.3$ 的指标对记为对立指标对。冲突因此同时在记录与指标对两个
层面度量，而不只由负相关定义。

\textbf{结果。}在 A1 全体上，$\tau=0.9$ 时的冲突率为 32.8\%，记录内指标的平均标准差为
0.260，对立指标对共 31 对。冲突在领域间差异很大：$\tau=0.9$ 时 arXiv 与 Book 的冲突率
分别为 97.8\% 与 98.2\%，而 StackExchange 与 CommonCrawl 仅为 13.2\% 与 18.8\%。

\textbf{成因分析。}六个打分体系针对的构念不同：规则类指标度量文本表面的格式与统计特征，
分类器类指标度量教育价值、流畅度、推理性等语义属性。冲突率在长文本、强格式的 arXiv 与
Book 中最高，与“不同构念在这类文本上给出相反判断”的解释一致；但这一解释来自冲突在领域间
的分布模式，并未对具体指标对逐一检验，本文不把它作为已证实的因果结论。

\textbf{消解规则。}式 \eqref{eq:q-record} 的体系平衡平均是本文的消解规则：同一体系内的
多个相关指标先合并为一票，再在体系之间等权，使冲突不会因某一体系指标数量多而被单方面
“投票”解决；冲突本身不被抹去，而是以冲突率与分歧幅度随评分一同报告。

\textbf{扩展集检验。}在 A2 与 A3 上，对立指标对的数量与抽样集完全相同（arXiv 为 26 对，
GitHub 为 31 对），冲突率曲线的最大变化不超过 0.005，主要结论在扩展集上同样成立。

\subsection{领域配比模型}

对第 $k$ 个带损失的领域（共 13 个），以配比为自变量建立岭回归\cite{hoerl1970ridge}
\begin{equation}
  L_k(\boldsymbol{p})=\theta_{k0}+\sum_{d\ne r}\theta_{kd}\,p_d
  +\sum_{d<e}\phi_{k,de}\,p_d\,p_e ,
  \label{eq:mixture}
\end{equation}
其中参照域 $r$ 取 Pile-CC（在 A4 上非零出现率与份额标准差之积最大），单纯形约束通过删去
参照域的一次项实现。候选模型为仅含一次项的线性模型（16 个特征）与加入全部两两交互项的
模型（152 个特征），岭参数按目标在 A4/A5 内做 5 折交叉验证选取。

模型选择规则在拟合前固定：只有当交互模型的交叉验证误差相对下降至少 5\%、且在至少 4 折中
占优时才采用。实际相对下降为 12.5\%，5 折全部占优，故采用交互模型。选择只使用 A4/A5，
模型在任何检验数据参与前即已冻结。

\begin{table}[htbp]
  \centering
  \caption{配比模型按检验角色的验证结果（13 个目标的宏平均）}
  \label{tab:q1-validation}
  \begin{tabular}{llrrrr}
    \toprule
    检验集 & 检验的命题 & RMSE & $R^2$ & Spearman & 平均相对误差 \\
    \midrule
    A6/A7 & 同规模（1M）留出集 & 0.394 & 0.725 & 0.872 & 7.2\% \\
    A8/A9 & 同一设计在 60M 的重复 & 1.521 & $-7.51$ & 0.875 & 43.2\% \\
    A10/A11 & 1B 规模、设计不同 & 2.896 & $-620.4$ & 0.849 & 145.4\% \\
    \bottomrule
  \end{tabular}
\end{table}

表\ref{tab:q1-validation}显示，模型在同规模留出集上通过检验。在 60M 上绝对预测失败：
各领域损失整体下降 0.97--2.44，但配对设计的排序几乎不变（逐目标 Spearman 为 0.980--0.998），
去掉各目标的均值后 $R^2$ 为 0.632，即\textbf{响应面的形状基本迁移而水平不迁移}。在 1B 上
规模与设计同时改变，去均值后的归一化 RMSE 为 2.208，超过预先设定的上限 2.0，形状也未能迁移。
因此该响应面只在 1M 规模上用于绝对预测，不作跨规模的绝对使用。

\textbf{领域效应。}在 A4 配比中心把 0.01 的份额由 Pile-CC 转给某一领域，该领域自身的验证
损失在全部 12 个可计算的情形中都下降，降幅为 0.0305（Wikipedia）至 0.1686（DM Mathematics）；
196 个跨领域效应中有 85 个的条件区间不含 0，存在替代与互补并存的结构。这些效应以冻结的模型
形式和惩罚参数为条件，只描述 1M 规模的响应面，不是领域质量的因果效应。

\subsection{质量评分与配比的结合}

A16 给出的参考映射只能把 17 个配比领域中的 6 个关联到质量领域，其余 11 个没有质量信息，
本文不为其插补任何评分。定义配比的已映射质量占比 $m(\boldsymbol{p})=\sum_{d\in\mathcal{M}}p_d$
与已映射部分的条件平均质量
\begin{equation}
  Q(\boldsymbol{p})=\frac{\sum_{d\in\mathcal{M}}p_d\,Q_d}{m(\boldsymbol{p})},\qquad m(\boldsymbol{p})>0,
\end{equation}
$m(\boldsymbol{p})=0$ 时 $Q(\boldsymbol{p})$ 无定义。在 A4 的 512 个设计上，$m$ 的平均值为
0.5646，这是已映射领域在配比中的实际份额，而不是已映射领域个数与 17 之比。在 1M 响应面上，$Q(\boldsymbol{p})$
平均只解释各目标 6.8\% 的变化，且只有 3 个目标的斜率为负。因此质量评分不是配比的充分概括，
本文不把 $Q$ 作为配比模型的输入，两者作为互补的信息分别报告。

\subsection{外推稳健性}

A13 与 A15 是估算参照，不是 10B 与 70B 规模上的观测。把 1M 模型直接用于这两张表时，
13 个目标的平均相对误差分别为 202\% 与 286\%，只能作为压力诊断。结合 60M 与 1B 的检验结果，本文
不给出 10B/70B 的配比外推结论；能够跨规模保留的只有同一设计下配比排序的稳定性，而不是
损失的绝对水平。
